\documentclass[11pt]{ctexart}
\usepackage{amsmath,amssymb,longtable}
\begin{document}
\section*{P416: 当前完整数学模型（2026-10-20）}
\section*{修订依据与边界}保留24货列及全部4客列，冲突客列P2整体顺延45分钟并留原时刻；有限路线按出现次数展开成状态无环图。完整分钟时刻/所有同反向安全间隔模型显式允许此有限路线扩展。

\section*{有限路线、分钟时刻与安全}
每个货列$j$有有限允许路线$r$（重复站点按出现次数展开），二元$y_{jr}$，$\sum_ry_{jr}\le1$。
每条候选路线的每次区间通行$e$有整数发到时刻$d_{jre},a_{jre}\in[0,H]$，$H=1440$。
\[
0\le d_{jre}\le Hy_{jr},\quad a_{jre}=d_{jre}+t_e y_{jr}.
\]
首发$d_{jr1}$在$[\max(0,I_j-S_j)y_{jr},\min(H,I_j+S_j)y_{jr}]$内。
偏移$\delta_{jr}\ge\pm(d_{jr1}-I_jy_{jr})$，$0\le\delta_{jr}\le S_jy_{jr}$。
中间站停留整数$w_e\in[m_ey_{jr},(m_e+\bar w_e)y_{jr}]$，
$d_{jre}=a_{jr,e-1}+w_e$。
同物理轨道的不同列车通行，用一个二元变量选择前后；若同向，前后出发至少差$g_d$、
到达至少差$g_a$且先后方向一致；若反向，后出发至少为前到达加$g_o$。
只有两条路线均选择时激活。所有固定客列逐区间时刻作为常数，实施同样互斥。
\[
\max\sum_{jr}\left[
\Pi_jy_{jr}-c^\delta_j\delta_{jr}
-c^w_j\sum_{e>1}(w_e-m_ey_{jr})
-c^r_j(\sum_{e\in r}t_e-T_j^{ideal})y_{jr}\right].
\]
所有时长为原边数据之和；不采用错误汇总值。P2整列时刻平移45分钟是输入冲突修复且留存原值。
循环物理路线扩展为有限状态路径，完整保留全部24个货列请求和4个固定客列；
站点可行性由上述容量和调度条件确定。

\end{document}
